

\chapter*{Conclusion}
\addcontentsline{toc}{chapter}{Conclusion}

Le rayonnement thermique constitue un mécanisme de transfert de chaleur incontournable dès lors que les températures deviennent élevées, que ce soit en régime nominal dans les réacteurs à haute température de type HTGR, ou en situation accidentelle dans les réacteurs à eau légère. La simulation de ces échanges radiatifs, en particulier dans des géométries tridimensionnelles complexes, constitue ainsi un enjeu important pour la plateforme multiphysique \trust du CEA.

Ce stage s'est inscrit dans cette problématique, avec pour objectif de consolider, d'étendre et de vérifier un module de calcul radiatif par lancer de rayons Monte Carlo (\emb), développé à partir d'une première maquette héritée d'un précédent stage. Ce travail a couvert l'ensemble de la chaîne de développement d'un tel outil : l'écriture du module en C++, avec une attention particulière portée aux problématiques de calcul haute performance (parallélisation, vectorisation), la mise en place d'une démarche de vérification rigoureuse du code au moyen de solutions analytiques et de comparaisons à un code de référence, puis la mise en œuvre de l'ensemble de la chaîne de simulation sur un cas d'étude représentatif, un assemblage de crayons combustible couplant calcul radiatif et simulation thermohydraulique 3D au sein de TRUST.

Cette dernière étape a nécessité la mise en commun de l'ensemble des briques développées au cours du stage : la prise en main de la plateforme \salome pour la génération des maillages surfaciques et volumiques, l'utilisation de TRUST pour la résolution du problème thermohydraulique couplé, ainsi que le développement de nombreux utilitaires Python dédiés au post-traitement des résultats, à la génération automatisée des maillages et à la création des fichiers d'entrée TRUST. Cette diversité d'outils, a ainsi pu être mobilisée de façon cohérente pour aboutir aux résultats du cas d'étude final.

Sur le plan personnel, ce stage a été l'occasion d'un apprentissage significatif du langage C++, jusqu'alors peu maîtrisé, appliqué directement à un contexte de calcul scientifique exigeant en performance. Il a également permis d'approfondir mes connaissances en transfert radiatif, en particulier dans le cas des milieux transparents entre surfaces opaques, ainsi que ma compréhension des enjeux propres à la vérification et à la validation d'un code de simulation. Enfin, la prise en main de deux outils centraux de l'écosystème logiciel du CEA, \trust et \salome, constitue un acquis qui dépasse le seul cadre de ce stage.

Ce travail ouvre plusieurs perspectives. Sur le plan fonctionnel, il ouvre la voie au développement de modèles à zones en trois dimensions, capables de tirer parti du module radiatif développé pour dépasser les limites des approches actuelles. Une telle démarche nécessiterait cependant d'explorer des méthodes alternatives au modèle des facteurs de vue d'ordre 1 utilisé dans ce rapport, afin de s'affranchir de l'hypothèse de radiosité uniforme dont l'étude du chapitre précédent a montré les limites. Plus largement, ce stage aura permis de partir d'une preuve de concept portant sur un modèle de rayonnement 3D pour aboutir à une preuve de concept de calcul CFD couplé, mettant en évidence la viabilité de cette chaîne de simulation pour des cas d'usage représentatifs. Enfin, ce travail mené en milieu transparent constitue également un premier pas vers ma thèse, intitulée \textit{Transferts de chaleur par rayonnement : résolution numérique efficace de problèmes associés en milieu beerien ou non pour les besoins de validation de modèles simplifiés}, qui prolongera cette démarche à des milieux participants.